% ══════════════════════════════════════════════════════════════════════════════
%  EXPERIMENT 2 — IMPACT CRATERS                         (Lesson 5, T4 Week 7)
% ══════════════════════════════════════════════════════════════════════════════
\experimentdiv{Impact Craters}{Dropping a model meteorite \textbullet{} Averages and anomalies \textbullet{} Reading the surface of the Moon}{5}\label{sec:2}
\lessoncodes{VC2S8U12, VC2S8H02 \textbullet{} VC2S8I01, I02, I04, I05, I06}

\vocabhead
\begin{ctab}
\begin{tabular}{|SW{4.2cm}|W{11.6cm}|}
  \hline\hd & \thead{Word} & \thead{Definition (copy from the slide)} \\ \hline
  & \scaf{\textbf{Impact crater}} & \scaf{A bowl-shaped hollow made when a meteorite or asteroid hits a planet or moon.} \\ \hline
  & \textbf{Meteorite} & \\ \hline
  & \textbf{Ejecta} & \\ \hline
  & \textbf{Diameter} & \\ \hline
  & \textbf{Anomaly} & \\ \hline
\end{tabular}
\end{ctab}

\begin{lipanel}
\inquiryli{Measure crater diameters precisely, average repeat drops, spot anomalies and draw a line of best fit.}
\understandli{Explain how the speed of an impact affects the size of a crater, and why the Moon has so many more craters than Earth.}
\sctitle
\begin{itemize}[leftmargin=1.4em, itemsep=0pt, topsep=1pt]
  \item I can measure a crater's diameter to the nearest millimetre.
  \item I can find an anomaly and leave it out of my average.
  \item I can explain my results using the speed of the falling marble.
\end{itemize}
\end{lipanel}

\begin{skillsfocus}
\focus{QP1}{Write a testable hypothesis.}
\focus{PC2}{Change only the drop height.}
\focus{PC4}{Measure each crater to the nearest mm.}
\focus{PA2}{Draw a line of best fit.}
\focus{PA3}{Calculate averages.}
\focus{PA7}{Find and explain anomalies.}
\focus{EV3}{Suggest how to reduce errors.}
\end{skillsfocus}

% ─── TECHNIQUE ────────────────────────────────────────────────────────────────
\exppart{Technique spotlight}{Blue}

\begin{technique}{Measuring a crater}{PC4}
\begin{minipage}[c]{0.6\linewidth}
Lift the marble out gently with tweezers. Lay a ruler across the crater, without touching the flour, and measure from rim to rim across the widest part. The \textbf{rim} is the raised edge, not the spray of cocoa (the \textbf{ejecta}). Read to the nearest millimetre. Look straight down at the ruler, not from the side.
\end{minipage}\hfill
\begin{minipage}[c]{0.36\linewidth}\centering
\begin{tikzpicture}[scale=0.75]
  \fill[black!8] (-2,-0.5) rectangle (2,0);
  \draw[Ink, line width=1pt] (-2,0) -- (-0.9,0) .. controls (-0.75,0.15) .. (-0.6,0.05) .. controls (-0.3,-0.45) and (0.3,-0.45) .. (0.6,0.05) .. controls (0.75,0.15) .. (0.9,0) -- (2,0);
  \foreach \x in {-1.7,-1.4,-1.2,1.15,1.4,1.75}{\fill[Ink] (\x,0.04) circle (0.04);}
  \draw[Ink, {Stealth[length=3pt]}-{Stealth[length=3pt]}] (-0.72,0.45) -- (0.72,0.45);
  \node[font=\scriptsize] at (0,0.75) {diameter};
  \node[font=\scriptsize, anchor=north] at (1.5,-0.55) {ejecta};
\end{tikzpicture}
\end{minipage}
\end{technique}

\begin{technique}{Anomalies}{PA7}
An \textbf{anomaly} is a result that does not fit the pattern of the others, such as 38, 41 and \textbf{57}\,mm from the same height. Circle it, suggest why it happened (the marble bounced, or hit a crater from an earlier drop), and leave it out of the average. Never just rub it out: a scientist records everything.
\end{technique}

\begin{workedex}{Averaging with an anomaly}
Drops from 50\,cm: 38, 41 and 57\,mm. 57 is an anomaly. \quad Average $= (38 + 41) \div 2 = 39.5$\,mm.
\end{workedex}

\begin{wedo}{Averaging with an anomaly}
Drops from 75\,cm: 46, 31 and 48\,mm. Find the anomaly, suggest a cause and calculate the average.
\workbox{3cm}
\end{wedo}

% ─── PLAN AND SAFETY / METHOD ────────────────────────────────────────────────
\exppart{Plan, safety and method}{Teal}

\textbf{Aim:} To investigate the effect of changing the drop height of a marble on the diameter of the crater it makes.

\task Write a hypothesis. Add a reason if you can. \skr{QP1}
\answerbox{3cm}

\needspace{12\baselineskip}
\task Complete the variables table. \skr{PC2}
\begin{ctab}
\begin{tabular}{|SW{4.4cm}|W{10.2cm}|}
  \hline\hd & \thead{Variable} & \thead{In this experiment} \\ \hline
  & \scaf{Independent} & \scaf{Drop height (cm)} \\ \hline
  & Dependent & \\ \hline
  & Controlled & \\ \hline
  & Controlled & \\ \hline
\end{tabular}
\end{ctab}

\subsubsection*{Equipment}
Deep tray of plain flour (at least 3\,cm deep) with a thin dusting of cocoa on top \textbullet{} one large marble \textbullet{} metre ruler \textbullet{} 30\,cm ruler \textbullet{} tweezers \textbullet{} sieve to smooth the flour \textbullet{} newspaper under the tray

\needspace{10\baselineskip}
\task Complete the risk assessment. \skr{PC3}
\begin{ctab}
\begin{tabular}{|SW{4.0cm}|W{5.2cm}|W{5.4cm}|}
  \hline\hd & \thead{Risk (what could go wrong)} & \thead{Prevention} & \thead{Responding (if it happens)} \\ \hline
  & \scaf{Flour dust gets in eyes} & \scaf{Drop gently; safety glasses; keep faces back from the tray} & \scaf{Rinse the eye with water; tell the teacher} \\ \hline
  & A marble on the floor is slipped on & & \\ \hline
  & Someone has a wheat or cocoa allergy & & \\ \hline
\end{tabular}
\end{ctab}

\subsubsection*{Method}
\begin{enumerate}[itemsep=2pt]
  \item Smooth the flour flat and sprinkle a thin layer of cocoa over it through the sieve.
  \item Hold the metre ruler upright beside the tray. Hold the marble so its bottom is 25\,cm above the flour, and let go without pushing.
  \item Measure the crater's diameter (rim to rim) to the nearest mm. Record it in Table 2.1.
  \item Drop twice more from 25\,cm, each time on an untouched part of the tray.
  \item Repeat steps 2 to 4 from 50, 75 and 100\,cm. Smooth the tray and add cocoa whenever it runs out of untouched space.
\end{enumerate}

% ─── RESULTS ─────────────────────────────────────────────────────────────────
\exppart{Results}{Amber}

\subsubsection*{Table 2.1: Crater diameter at each drop height \skr{PA1}}
\begin{ctab}
\begin{tabular}{|SW{3.0cm}|C{2.7cm}|C{2.7cm}|C{2.7cm}|C{3.0cm}|}
  \hline\hd & \thead{Drop height (cm)} & \thead{Drop 1 (mm)} & \thead{Drop 2 (mm)} & \thead{Drop 3 (mm)} & \thead{Average (mm)} \\ \hline
  & \scaf{Example: 10} & \scaf{29} & \scaf{31} & \scaf{30} & \scaf{30.0} \\ \hline
  & 25 & & & & \\ \hline
  & 50 & & & & \\ \hline
  & 75 & & & & \\ \hline
  & 100 & & & & \\ \hline
\end{tabular}
\end{ctab}

\task Circle any anomaly in Table 2.1, and show your working for one average. \skr{PA3}
\workbox{3cm}

\needspace{22\baselineskip}
\task Draw a line graph of the average diameter against drop height, with a title and labelled axes. Draw a smooth line of best fit. \skr{PA2}
\graphgrid{10cm}

% ─── ANALYSIS ────────────────────────────────────────────────────────────────
\exppart{Analysis and evaluation}{Coral}

\task Describe the trend. Use numbers from Table 2.1. \skr{PA5}
\answerbox{4cm}

\needspace{10\baselineskip}
\task Explain why a higher drop made a bigger crater. Use the idea of speed (and energy, if you can).
\answerbox{4cm}

\needspace{10\baselineskip}
\task Describe any anomaly you found and suggest what caused it. If you found none, say what could have caused one. \skr{PA7}
\answerbox{4cm}

\needspace{12\baselineskip}
\task Give two sources of error and how each could be reduced. \skr{EV3}
\begin{ctab}
\begin{tabular}{|SW{7.2cm}|W{7.8cm}|}
  \hline\hd & \thead{Source of error} & \thead{How to reduce it} \\ \hline
  & \scaf{The marble was pushed, not dropped} & \scaf{Hold it on a spoon and tip it off} \\ \hline
  & & \\ \hline
  & & \\ \hline
\end{tabular}
\end{ctab}

\needspace{12\baselineskip}
\task The Moon is covered in craters but Earth has very few. Give two reasons. Think about the air, the weather and plate movement.
\answerbox{4.5cm}

\needspace{10\baselineskip}
\task Write a conclusion: the trend, the numbers that back it, and whether your hypothesis was supported. \skr{EV1}
\answerbox{4.5cm}

\begin{yr8box}
\begin{keepext}
\exttask Tnorala (Gosses Bluff) in the Northern Territory is the remains of a crater about 22\,km across, made around 140 million years ago. The Western Arrernte account of Tnorala tells of a baby's wooden carrier that fell from the sky and struck the Earth. Explain how this account and the scientific explanation agree, and plan an experiment with your tray to test whether a heavier meteorite makes a wider crater.
\answerbox{6cm}
\end{keepext}
\end{yr8box}

% ─── SKILL CHECK ─────────────────────────────────────────────────────────────
\exppart{Skill check}{NavyBlue}

\begin{skillcheck}
\scrow{QP1}{I can write a hypothesis that names both variables.}
\scrow{PC2}{I changed only the drop height.}
\scrow{PC4}{I measured each crater rim to rim, to the nearest mm.}
\scrow{PA2}{I drew a smooth line of best fit.}
\scrow{PA3}{I calculated averages, leaving out anomalies.}
\scrow{PA7}{I can spot an anomaly and suggest its cause.}
\scrow{EV3}{I can suggest ways to reduce errors.}
\end{skillcheck}

\task Look back over both experiments and the whole year. Which science skill have you improved most in Year 7? Give an example.
\answerbox{5cm}

\end{document}
